\section{Sets with Few Distances}

The unit distance problem fixes the points and counts the distances between them. We now turn this around, and ask how many points there can be if they determine only one distance, or only two.

\subsection{Equilateral Sets}

Let's now ask ourselves a question that, though seemingly similar to the unit distance problem, is actually quite different:
\begin{center}
    How many points in $\R^d$ can all be at pairwise unit distance?
\end{center}
In the unit distance problem, we \textit{fixed points} and \textit{counted unit distances}. Here, we haven't fixed any points.

It turns out the answer is $d + 1$, and is given by the regular simplex.

Let's study the properties of hypothetical points that would satisfy this condition to glean information that will help us figure out what they might be.

% [CORRECTED] was: "... with $i \neq j$ and $\norm{v_i} = 1$ for $i \neq 0$. Without loss of generality, take $v_0 = 0$." --- assumed before the translation, $\norm{v_i} = 1$ is an extra hypothesis (the origin at distance $1$ from the points); after it, it follows from $\norm{v_i - v_0} = 1$
Let $v_0, v_1, \ldots, v_n \in \R^d$ be points such that $\norm{v_i - v_j} = 1$ for all $0 \leq i, j \leq n$ with $i \neq j$. Without loss of generality, take $v_0 = 0$, so that $\norm{v_i} = 1$ for $i \neq 0$.

For all $i, j \geq 1$ with $i \neq j$, notice that
\begin{align*}
    1 = \norm{v_i - v_j}^2 = \cycl{v_i - v_j, v_i - v_j} = \cycl{v_i, v_i} - 2\cycl{v_i, v_j} + \cycl{v_j, v_j} = 2\parenth{1 - \cycl{v_i, v_j}}
\end{align*}
This tells us
\begin{align*}
    \cycl{v_i, v_j} =
    \begin{cases}
        1 & \text{ if } i = j \\
        \frac{1}{2} & \text{ if } i \neq j
    \end{cases}
\end{align*}

\begin{claim*}
    $v_1, \ldots, v_n$ are linearly independent.
\end{claim*}
\begin{proof}
    \sorry % [CLAUDE] leave this, this one is very easy
\end{proof}

\sorry % the lecture stops at the claim: the bound on $n$, and so on the number of points, is not yet drawn from it

\subsection{Two-Distance Sets}

What if I have two distances instead?

Suppose $p_1, \ldots, p_n \in \R^d$ are at pairwise distances $\set{\alpha, \beta}$. We want to define polynomials $f_i(\bx)$ such that for all $i, j$,
\begin{align*}
    f_i\of{p_j} = 
    \begin{cases}
        \alpha^2 \beta^2 & \text{ if } i = j \\
        0 & \text{ if } i \neq j
    \end{cases}
\end{align*}
We can define
\begin{align*}
    f_i(\bx) = \parenth{\norm{\bx - p_i}^2 - \alpha^2}\parenth{\norm{\bx - p_i}^2 - \beta^2}
\end{align*}

\sorry % the lecture breaks off here, and the author has asked about it for office hours; read the issue below before filling this
% Question for office hours: https://github.com/thefundamentaltheor3m/DiscreteMathNotes/issues/33 (the two-distance case)
